\section*{Appendix}

\tcbset{
  evoCase/.style={
    enhanced,
    breakable,
    arc=1mm,
    boxrule=0.45pt,
    left=4pt,
    right=4pt,
    top=4pt,
    bottom=4pt,
    fonttitle=\bfseries\small,
    coltitle=black
  },
  skillCard/.style={
    enhanced,
    arc=1mm,
    boxrule=0.35pt,
    left=4pt,
    right=4pt,
    top=3pt,
    bottom=3pt,
    fonttitle=\bfseries\footnotesize,
    coltitle=black
  }
}

\section{Benchmark Details}\label{app:benchmarks}

Table~\ref{tab:bench-stats} summarizes the benchmark settings. We use the same
online protocol across benchmarks: tasks are processed as a stream, the solver
uses the current harness, and the evolver updates the harness after observing
task outcomes.

\begin{table}[h]
\centering
\caption{Benchmark statistics and evaluation signals.}
\label{tab:bench-stats}
\small
\setlength{\tabcolsep}{3pt}
\scalebox{0.78}{
\begin{tabular}{lrll}
\toprule
\textbf{Benchmark} & \textbf{\#Tasks} & \textbf{Domain} & \textbf{Evaluation signal} \\
\midrule
WebArena Inf. & 80 & Web navigation & State verifier \\
Terminal-Bench & 89 & CLI / scripting & Docker verifier \\
SWE-bench Lite & 300 & Software eng. & Unit tests \\
CL-Bench & 1899 & Adaptive reasoning & Rubric judge \\
TAU-Bench & 165 & Tool use & State verifier \\
\bottomrule
\end{tabular}}
\end{table}

\paragraph{WebArena Infinity.}
We use 80 hard tasks from the email-centered ``superhuman-general'' setting and
related web applications. Tasks require stateful browser interaction, such as
creating labels, changing settings, using email filters, or coordinating
calendar/email information. A programmatic verifier checks the final application
state. This benchmark is useful for studying whether skills capture UI
procedures.

\paragraph{Terminal-Bench v2.}
We evaluate 89 tasks across software engineering, data processing, scientific
computing, security, system administration, and related command-line domains.
The solver works in an isolated shell environment and submits final artifacts or
answers inside a task-specific Docker image. Final success is determined by the
benchmark's executable verifier, typically a hidden or mounted test script run
inside the container. 

\paragraph{SWE-bench Lite.}
We use 300 real GitHub issues from 12 Python repositories. The solver edits code
and the final patch is evaluated by repository tests.

\paragraph{CL-Bench.}
CL-Bench contains 1899 adaptive reasoning tasks across Domain Knowledge
Reasoning, Empirical Discovery \& Simulation, Procedural Task Execution, and
Rule System Application. Each task is graded with rubric criteria.

\paragraph{TAU-Bench.}
TAU-Bench contains 165 customer-service tool-use tasks across airline and
retail domains. Evaluation checks the final database state.

\section{Harness Artifacts Examples}\label{app:examples}

This section shows representative learned artifacts without enumerating the
learned inventory. For task-level comparison, we align the same task identifiers
between the no-evolve and evolved runs. We only call a case a direct improvement
when the baseline run fails and the evolved run succeeds. This is not a
single-skill causal ablation: several skills may be injected together, and some
benchmarks do not record injected skills in every result file. The examples
therefore use a structured failure/success format: task, verifier signal,
trajectory summary, learned skill content, and a short takeaway. 

\begin{tcolorbox}[evoCase, colback=blue!2, colframe=blue!45,
title={WebArena: verified setting}]
\small
\begin{itemize}[leftmargin=1.1em, itemsep=1pt, topsep=1pt]
\item \textbf{Task:} \texttt{task\_h10}; create ``Legal'' and disable
``Support Ticket''.
\item \textbf{Baseline:} verifier found ``Support Ticket'' still enabled,
despite the trajectory claiming both subgoals were complete.
\item \textbf{Evolved:} verifier passed; the trajectory returned to the Auto
Labels list and checked the existing label state.
\item \textbf{Learned skill:} toggle an existing auto label off; do not delete
it; verify creation and disabling independently.
\end{itemize}
\tcblower
\textbf{Core point.} The harness learned a concrete UI distinction:
create-new-label and disable-existing-label are separate operations.
\end{tcolorbox}

\begin{tcolorbox}[evoCase, colback=blue!2, colframe=blue!45,
title={WebArena: grounded recipients}]
\small
\begin{itemize}[leftmargin=1.1em, itemsep=1pt, topsep=1pt]
\item \textbf{Task:} \texttt{task\_h65}; reply to Kevin Zhao with CTO and VP
Engineering on CC.
\item \textbf{Baseline:} sent to plausible names but wrong domains
(\texttt{acme.com} instead of \texttt{acmecorp.com}).
\item \textbf{Evolved:} verifier passed with the required role-derived
recipients in CC.
\item \textbf{Learned skill:} prefer authoritative role sources and preserve
exact email addresses across apps.
\end{itemize}
\tcblower
\textbf{Core point.} The harness learned identity grounding, not merely email
composition.
\end{tcolorbox}

\begin{tcolorbox}[evoCase, colback=orange!2, colframe=orange!45,
title={Terminal-Bench: executable recovery}]
\small
\begin{itemize}[leftmargin=1.1em, itemsep=1pt, topsep=1pt]
\item \textbf{Task:} \texttt{adaptive-rejection-sampler}; pass the Docker
verifier for an R sampler.
\item \textbf{Baseline:} formal tests failed after repeated reruns of the same
R test command.
\item \textbf{Evolved:} passed after diagnosing timeout, missing dependency,
formal-test, and sorted-vector failures.
\item \textbf{Learned skill:} after repeated similar failures, stop, shrink the
case, build a minimal reproducer, then add complexity back.
\end{itemize}
\tcblower
\textbf{Core point.} The harness learned a recovery policy for executable
feedback loops.
\end{tcolorbox}

\begin{tcolorbox}[evoCase, colback=green!2, colframe=green!45,
title={SWE-bench: test-grounded repair}]
\small
\begin{itemize}[leftmargin=1.1em, itemsep=1pt, topsep=1pt]
\item \textbf{Task:} \texttt{django\_\_django-11133}; fix Django response
handling without regressing existing behavior.
\item \textbf{Baseline:} timed out after source search and failed test attempts.
\item \textbf{Evolved:} passed \texttt{FAIL\_TO\_PASS: 1/1} and
\texttt{PASS\_TO\_PASS: 64/64}.
\item \textbf{Learned skill:} read complete failing-test assertions, match
exact expected behavior, then run targeted tests before submission.
\end{itemize}
\tcblower
\textbf{Core point.} The harness learned to anchor patches in executable test
expectations.
\end{tcolorbox}

\begin{tcolorbox}[evoCase, colback=purple!2, colframe=purple!45,
title={TAU-Bench: complete tool execution}]
\small
\begin{itemize}[leftmargin=1.1em, itemsep=1pt, topsep=1pt]
\item \textbf{Task:} \texttt{retail\_task\_5}; exchange multiple items from one
order.
\item \textbf{Baseline:} database verifier failed, consistent with partial or
wrong-tool execution.
\item \textbf{Evolved:} verifier passed after retrieving the order and applying
all requested item changes.
\item \textbf{Learned skill:} enumerate every requested operation, choose the
tool by order status, and verify the response covers the full list.
\end{itemize}
\tcblower
\textbf{Core point.} The harness learned complete tool-call execution.
\end{tcolorbox}

\section{Which Tasks Are Easier to Improve?}\label{app:improvement-hotspots}

The most informative pattern is not the total number of fail-to-pass cases, but
the kind of task that turns around after harness evolution.
Table~\ref{tab:improvement-hotspots} groups matched fail-to-pass cases by their
observable failure mode and gives representative task identifiers from the
runs in \texttt{outputs}.

\begin{table*}[t]
\centering
\caption{Improvement hotspots from matched baseline-fail/evolved-pass cases.}
\label{tab:improvement-hotspots}
\small
\setlength{\tabcolsep}{3.5pt}
\begin{tabular}{p{2.0cm}p{3.2cm}p{4.3cm}p{4.0cm}}
\toprule
\textbf{Benchmark} & \textbf{Easier-to-improve tasks} &
\textbf{Representative cases} & \textbf{Why skills help} \\
\midrule
WebArena &
Stateful UI tasks with exact recipients, labels, settings, and mail state. &
\texttt{h10}: disable an existing auto label;
\texttt{h65}: preserve corporate email domains;
\texttt{h73}: archive all qualifying read-important emails. &
The verifier checks final app state, so procedural memory about where to click
must be coupled with exact state confirmation. \\
\midrule
Terminal-Bench &
Tasks with executable feedback loops and recoverable failures. &
\texttt{adaptive-rejection-sampler}, \texttt{dna-insert},
\texttt{build-pmars}, \texttt{regex-chess}. &
The solver can observe Docker-test failures; process skills help it pivot,
reduce scope, and verify artifacts before submission. \\
\midrule
SWE-bench Lite &
Mature repositories with stable conventions, especially Django and Sympy. &
\texttt{django\_\_django-11133}, \texttt{django\_\_django-17087},
\texttt{sympy\_\_sympy-19254}. &
The issue often hides exact expectations in tests or framework internals; topic
skills preserve these conventions across later tasks. \\
\midrule
TAU-Bench &
Retail exchanges/returns and airline reservation edits with multiple
constraints. &
\texttt{retail\_task\_5}, \texttt{retail\_task\_91},
\texttt{airline\_task\_13}. &
The state verifier rewards complete tool-call execution. Skills prevent
partial updates, wrong mutation tools, and guessed item or flight identifiers. \\
\midrule
CL-Bench &
Rule-application, procedural, and domain-knowledge contexts with local rubrics.
&
Context categories with many turnarounds include Rule System Application,
Procedural Task Execution, and Domain Knowledge Reasoning. &
Context skills behave like small local manuals: exact extraction rules, role
boundaries, output schemas, and calculation conventions. \\
\bottomrule
\end{tabular}
\end{table*}

\section{What the Harness Actually Learns}\label{app:learned-skills}

Figure~\ref{fig:learned-skill-excerpts} shows real skill excerpts selected with
matched fail-to-pass support. The evidence line reports the relevant turnaround
family in the runs, while the rule line gives the core learned behavior.

\begin{figure*}[t]
\centering
\begin{minipage}[t]{0.48\textwidth}
\begin{tcolorbox}[skillCard, colback=blue!2, colframe=blue!45,
title={WebArena: UI-state verification}]
\small
\textbf{Support.} Settings-state tasks account for 24\% of WebArena fail-to-pass cases; related UI skills appear in 72\% of improved runs.\\
\textbf{Guidance.} Treat creation and disabling as separate operations; after editing an auto label, return to the settings list and verify its final state.
\end{tcolorbox}
\end{minipage}
\hfill
\begin{minipage}[t]{0.48\textwidth}
\begin{tcolorbox}[skillCard, colback=blue!2, colframe=blue!45,
title={WebArena: grounded identity use}]
\small
\textbf{Support.} Identity-grounding failures account for 32\% of WebArena fail-to-pass cases.\\
\textbf{Guidance.} Use authoritative sources for role contacts and preserve exact email addresses; do not infer or normalize domains.
\end{tcolorbox}
\end{minipage}

\vspace{4pt}
\begin{minipage}[t]{0.48\textwidth}
\begin{tcolorbox}[skillCard, colback=orange!2, colframe=orange!45,
title={Terminal-Bench: executable recovery}]
\small
\textbf{Support.} Executable-feedback and recovery patterns account for 100\% of Terminal-Bench fail-to-pass cases.\\
\textbf{Guidance.} Stop rerunning the same command, isolate a minimal failing case, fix the core issue, and then rebuild toward the full solution.
\end{tcolorbox}
\end{minipage}
\hfill
\begin{minipage}[t]{0.48\textwidth}
\begin{tcolorbox}[skillCard, colback=green!2, colframe=green!45,
title={SWE-bench: test-grounded repair}]
\small
\textbf{Support.} Django and Sympy account for 66.7\% of SWE-bench fail-to-pass cases.\\
\textbf{Guidance.} Read the full failing-test assertions before editing; match exact expected strings, error identifiers, and output structures.
\end{tcolorbox}
\end{minipage}

\vspace{4pt}
\begin{minipage}[t]{0.48\textwidth}
\begin{tcolorbox}[skillCard, colback=purple!2, colframe=purple!45,
title={TAU-Bench: complete tool execution}]
\small
\textbf{Support.} Retail workflows account for 85.2\% of TAU-Bench fail-to-pass cases; this guidance appears in 100\% of improved runs.\\
\textbf{Guidance.} Enumerate every requested operation, choose tools according to the current state, and verify that the response covers the full request.
\end{tcolorbox}
\end{minipage}
\hfill
\begin{minipage}[t]{0.48\textwidth}
\begin{tcolorbox}[skillCard, colback=gray!5, colframe=gray!45,
title={CL-Bench: constrained extraction}]
\small
\textbf{Support.} Rule, procedural, and domain-knowledge contexts dominate CL-Bench fail-to-pass cases.\\
\textbf{Guidance.} For extraction tasks, return the requested source text verbatim and match only the exact term or field specified by the prompt.
\end{tcolorbox}
\end{minipage}

\caption{
Representative learned harness excerpts across benchmarks. Each card pairs a recurring fail-to-pass pattern with the corresponding evolved guidance. Percentages summarize matched baseline-fail/evolved-pass cases and contextualize the examples rather than establish single-skill causality.
}
\label{fig:learned-skill-excerpts}
\end{figure*}

\section{Prompt Templates}\label{app:prompts}

The prompts follow the same context-to-harness compilation logic across experiments: the solver reflects on an execution context to propose candidate lessons, and the evolver curates these candidates against the existing harness. Below, we provide representative prompt templates that capture this common proposal-and-curation process.

\subsection{Solver-Side Proposal Prompt}

\begin{tcolorbox}[breakable, colback=gray!5, colframe=gray!50, fontupper=\small,
title={\small Proposal prompt, compressed}]
\textbf{Inputs:} evaluation result, verifier details or rubric feedback,
trajectory signals, compressed trajectory, and related existing skills.\\[2pt]
\textbf{Analyze the execution.}
For each distinct issue, identify what went wrong and the concrete missing
action, command, tool call, navigation step, or domain rule.\\[2pt]
\textbf{Propose a reusable skill.}
Choose a broad topic, decide \texttt{NEW}, \texttt{ENHANCE}, or \texttt{NONE},
write a trigger-style description, and provide short bullet-point content with
techniques and gotchas.\\[2pt]
\textbf{Filter aggressively.}
Skip generic advice, basic tool usage, exact task replay, and skills that would
not help unseen tasks.
\end{tcolorbox}

\subsection{Evolver-Side Curator Prompt}

\begin{tcolorbox}[breakable, colback=gray!5, colframe=gray!50, fontupper=\small,
title={\small Curator prompt, compressed}]
\textbf{Inputs:} topic, current skill library, budget, and proposals from the
latest batch.\\[2pt]
For each proposal, choose \texttt{ACCEPT}, \texttt{MERGE}, or \texttt{SKIP}.
Prefer merging over duplication, respect the budget, require a clear trigger
description, keep content short and actionable, and apply a generalizability
test such as usefulness for multiple unseen tasks.
\end{tcolorbox}

\subsection{General Skill Curator}

\begin{tcolorbox}[breakable, colback=gray!5, colframe=gray!50, fontupper=\small,
title={\small General skill curator, compressed}]
Analyze repeated patterns across tasks rather than within a single topic.
Create or update a general skill only when a pattern appears across multiple
contexts. General skills must avoid context-specific references and encode
procedures that can guide planning, verification, recovery, or tool use across
tasks.
\end{tcolorbox}

\section{Experimental Details}
\label{app:details}
All experiments use the same open-source pipeline with benchmark-specific adapters. WebArena uses browser workers with persistent application state. Terminal-Bench uses task-specific Docker images and executable verifiers. SWE-bench uses isolated repository environments and unit-test evaluation. TAU-Bench uses the official tool-use environment and verifier. CL-Bench uses a rubric-based judge pipeline.

The harness is stored as Markdown skill files with lightweight YAML metadata, making the learned artifact inspectable and easy to transfer across runs. Each entry contains a trigger describing when it should be retrieved, a short actionable rule or procedure, optional evidence linking it to prior executions, and a scope indicating whether it is intended as cross-task guidance or task-type guidance. For harness selection, we use Claude Sonnet 4.5 across all experiments to retrieve relevant skills from the current harness before task execution.

Unless otherwise specified, we use a batch size of 16. We set the maximum number of general skills and the maximum number of skills under each task-type topic to 5. This budget encourages the evolver to merge overlapping guidance and avoid accumulating overly detailed or task-specific records.

\section{Potential Risks}
\algname{} evolves skill libraries that are injected into agent prompts as procedural guidance. The learned skills are designed to capture task-solving patterns such as navigation flows, tool-use procedures, verification steps, and recovery strategies. In our experiments, all benchmarks are run in sandboxed environments with no access to real user data, external accounts, or production systems. Therefore, the evaluated setting poses minimal practical risk beyond standard benchmark execution.

\section{Use Or Create Scientific Artifacts}
Our work uses five established public benchmarks: WebArena-Infinity, TerminalBench2, SWE-bench Lite, CL-Bench, and TAU-Bench. We do not modify the benchmarks themselves. We use foundation models (Claude and other models) via their standard AWS Bedrock APIs. We will release the \algname{} pipeline code and all evolved skill libraries upon publication.

\subsection{Cite Creators Of Artifacts}
All benchmarks and models are properly cited in the main text.

\subsection{Discuss The License For Artifacts}
All benchmark resources used in this work are publicly available under their respective research or open-source licenses. In particular, WebArena Infinity and TAU-Bench are released under MIT licenses, WebArena is released under Apache-2.0, and CL-Bench is released under a custom evaluation-only license. Our pipeline code will be released under the MIT License.

\subsection{Artifact Use Consistent With Intended Use}
All benchmarks are used for their intended purpose of evaluating agent capabilities. The foundation models are used via their standard APIs in accordance with their terms of service.

\subsection{Data Contains Personally Identifying Info Or Offensive Content}
The benchmarks use synthetic or sanitized scenarios. SWE-bench uses public GitHub issues. To our knowledge, none contain personally identifying information or offensive content.

\section{Computational Experiments}\label{sec:app-compute}

\subsection{Compute Budget}
The full experimental suite is computationally expensive because it covers multiple benchmarks, solver models, evolution settings, and ablation configurations. Across all main experiments and diagnostic runs, the total API cost was on the order of \(\$100\mathrm{K}\). For a single benchmark configuration, we typically run 16 parallel workers. Depending on benchmark complexity, task timeout, and environment overhead, one full run takes approximately 2--8 hours of wall-clock time. Web and tool-use benchmarks are usually faster, while software-engineering and command-line tasks tend to take longer because they require repository setup, Docker execution, test running, or verifier calls.

\subsection{Experimental Setup And Hyper-params}
Key hyperparameters are described in Appendix~\ref{app:details}. We use fixed random seeds (42) for task shuffling to ensure reproducibility.

\subsection{Descriptive Statistics}
We report pass rate (\%) as the primary metric across all benchmarks. For statistical reliability, we report results over the full task set for each benchmark rather than subsampling. The performance is averaged over three runs.

\subsection{Parameters For Packages}

The project primarily used AWS Bedrock API for foundation model inference, together with standard AWS infrastructure services such as EC2, S3, EBS/FSx, CloudWatch, IAM, and VPC. The software environment used Python 3.10+, boto3 1.34+, botocore 1.34+, PyTorch 2.0+, CUDA 12.0+, Docker, and Ubuntu 24.04.

\section{AI Assistants In Research Or Writing}

This paper studies LLMs as the research object and uses frozen LLM backbones
for experimental inference. These models are not used to create new benchmark
labels or to fabricate experimental results. Reported numerical results come
from the experimental code and are checked by the authors.

AI assistant tools were used for polishing manuscript. They were not used to
decide scientific claims, select reported results, or replace author
verification.
